% ============================================================
%  Setmana 13 — Unitat 13 · Probabilitat
%  Exercicis: 28, 29, 40, 60, 62, 32, 33, 35
% ============================================================

\setmana{13}{Unitat 13 \textperiodcentered{} Probabilitat}{13 de desembre de 2026}%
        {28, 29, 40, 60, 62, 32, 33 i 35}

\begin{errada}{Exercici 40 a)}
  \ondiu{$VR_{5,6}=5^6=15\,625$}
  \hadedir{$4\cdot VR_{5,5}=4\cdot 5^5=12\,500$}
  \motiu{Els dígits parells són $0$, $2$, $4$, $6$ i $8$, però un nombre de
    6 xifres no pot començar per $0$: per a la primera xifra només hi ha
    4 opcions. El resultat $15\,625$ compta també seqüències com $002468$, que
    no són nombres de 6 xifres.}
\end{errada}

\begin{errada}{Exercici 60 a)}
  \ondiu{$P(A\cap B)=1-P(\overline{A}\cap\overline{B})=1-0{,}3=0{,}7$}
  \hadedir{$P(A\cup B)=1-P(\overline{A}\cap\overline{B})=1-0{,}3=0{,}7$}
  \motiu{«No s'ha complert alguna de les dues normes» és la unió $A\cup B$,
    que és el contrari de $\overline{A}\cap\overline{B}$ (complir-les totes
    dues). La intersecció val una altra cosa:
    $P(A\cap B)=0{,}23+0{,}65-0{,}7=0{,}18$.}
\end{errada}

%\senseerrades{28, 29, 62, 32, 33 i 35}

% ---------- Observacions (només surten amb \obstrue) ----------
\ifobs
\bigskip
\noindent
Una observació menor:\par\nopagebreak

\begin{observacio}{Exercici 35}
  A la solució hi ha una errada tipogràfica sense importància: hi diu
  «probabiltat» en lloc de «probabilitat».
\end{observacio}
\fi
